\documentclass{article}

\begin{document}
Description of simulation of conditioned S-I-R epidemic in population
of n individuals:

Unconditioned case: innovation proposals are updated by an MCMC chain
as follows.
\begin{itemize}
\item[A:] renew removals as Poisson process marked by uniformly random
  layers;
\item[A':] renew marks of removal proposals;
\item[B1-n:] iterate through layers 1 to n:
  \begin{itemize}
  \item renew infection proposals attached to a layer as Poisson
    process marked by subordinate layers chosen uniformly at random,
    as well as by the layers to which they are attached.
\end{itemize}
\end{itemize}

Step A' appears redundant! However, we need to impose conditioning via
the functional that converts the configuration of infection and
removal proposals into an epidemic trajectory defined as a sequence of
incidents. The conditioning is of constraint form: the epidemic
removal epochs are held fixed, prescribed by the conditioning.
Correspondingly, the MCMC chain rejects all proposals which violate
this constraint. Detailed balance arguments show the conditioned
epidemic is an equililibrium distribution for: the constrained MCMC
chain; the moves of A' now ensure that this new chain is irreducible
and therefore that the equilibrium is unique. (Otherwise an implicit
bias would be introduced via a subtle mechanism involving the fixed
marks of the conditioned removals.)

We require a further refinement in order to maintain a kind of
monotonicity (important for perfect simulation). Each process is
updated ``pixel-by-pixel''; in practice we generate the new process,
then consider additions of proposals from the new process and
deletions of proposals from the old process in time-order. We accept
each of the additions and deletions if as a result the conditioning is
not violated.
 
VOCABULARY: ACTIVATION AND  DECLINING  VERSUS ACCEPTANCE AND REJECTION 

To understand the conditioning, we must now describe the functional
which converts a set of innovation proposals into an epidemic
trajectory. The conversion procedure is as follows. First, set apart
layers corresponding to the initial infectives. Collect all infection
and removal proposals, and order them by time. Iterate as follows:
move to the next proposal for which the marks (one in case of removal,
two in case of infection) are layers corresponding to infectives.
Decline all previous proposals. Activate this proposal and repeat,
till either all proposals are exhausted or there are no more
infectives. ``Activation'' is as follows: if the activated proposal is a
removal then move one of the infective layers over to a ``removed''
category. If the activated proposal is an infection then add to the
infectives a layer which previously had not been among the infectives
(the collection of such layers forms the ``susceptible'' category!).

TO PROVE: RESULT IS INDEED AN S-I-R EPIDEMIC. 

Implementation details vary from this description, of course,
especially since for the conditional simulation there is a need for
each proposal to compute aspects of the trajectory in order to
determine whether the conditioning has been violated. we alleviate the
computational burden here by tracking for each layer: (a) the latest
declined infection proposal which if activated would lead to
activation of a removal which should be declined (the ``nearly bound'')
and (b) the latest activated proposal which if declined would lead to
the declining of a removal proposal whose activation is required in
order to satisfy the conditioning.

TO PROVE: CONSTRAINT LENDS TO DESIRED CONDITIONING 

TO PROVE: MONOTONICITY?  

\end{document}
